% =====================================================================
%  LiaScript / Markdown CheatSheet — A4, doppelseitig
%  Workshop "Lehre aktivierend gestalten mit LiaScript", 02.-09.06.2026
%  Compile:  pdflatex cheatsheet.tex   (zweimal für korrekte Boxen)
% =====================================================================
\documentclass[a4paper,landscape,9pt]{extarticle}

% --- Layout ----------------------------------------------------------
\usepackage[landscape,margin=0.9cm,top=1.0cm,bottom=0.7cm]{geometry}
\usepackage{multicol}
\setlength{\columnsep}{8pt}
\setlength{\parindent}{0pt}
\setlength{\parskip}{0pt}

% --- Sprache & Schrift ----------------------------------------------
\usepackage[T1]{fontenc}
\usepackage[utf8]{inputenc}
\usepackage[ngerman]{babel}
\usepackage{libertine}
\usepackage[varqu,varl,scaled=0.92]{inconsolata}

% --- Farben ----------------------------------------------------------
\usepackage{xcolor}
\definecolor{mdblue}{HTML}{1F4E79}
\definecolor{mdblueLt}{HTML}{E8EEF5}
\definecolor{liagreen}{HTML}{2E6B3D}
\definecolor{liagreenLt}{HTML}{E8F0E9}
\definecolor{codebg}{HTML}{F6F6F4}
\definecolor{rulegray}{HTML}{D0D0D0}
\definecolor{muted}{HTML}{606060}

% --- Pakete ----------------------------------------------------------
\usepackage[most]{tcolorbox}
\usepackage{listings}
\usepackage{qrcode}
\usepackage{fontawesome5}
\usepackage{array}
\usepackage[normalem]{ulem}
\usepackage{hyperref}
\hypersetup{hidelinks}

% --- Listings: Markdown/LiaScript-Quellcode --------------------------
\lstdefinelanguage{lia}{
  morecomment=[l]{\#},
  morecomment=[s]{<!--}{-->},
  sensitive=true
}
\lstset{
  basicstyle=\ttfamily\scriptsize,
  breaklines=true,
  breakatwhitespace=true,
  showstringspaces=false,
  keepspaces=true,
  columns=fullflexible,
  upquote=true,
  literate=%
    {ä}{{\"a}}1 {ö}{{\"o}}1 {ü}{{\"u}}1
    {Ä}{{\"A}}1 {Ö}{{\"O}}1 {Ü}{{\"U}}1
    {ß}{{\ss}}1 {„}{{\glqq}}1 {“}{{\grqq}}1,
  frame=none,
  xleftmargin=0pt,
  aboveskip=0pt,
  belowskip=0pt
}

% --- Karten-Vorlage --------------------------------------------------
%  Zwei Varianten: mdcard (blau, Markdown) und liacard (grün, LiaScript)
%  Innen: links Quellcode, rechts Wirkung — getrennt durch \tcblower.
\tcbset{
  cheatbase/.style={
    enhanced,
    sidebyside,
    sidebyside align=top,
    lefthand width=0.46\linewidth,
    segmentation style={solid, rulegray, line width=0.3pt},
    colback=white,
    boxrule=0.3pt,
    arc=2pt,
    left=4pt,right=4pt,top=2pt,bottom=2pt,middle=4pt,
    before skip=3pt, after skip=3pt,
    breakable=false,
    fonttitle=\sffamily\bfseries\footnotesize,
    coltitle=white
  }
}

\newtcolorbox{mdcard}[1]{cheatbase, colframe=mdblue!75,  colbacktitle=mdblue,   title=#1}
\newtcolorbox{liacard}[1]{cheatbase, colframe=liagreen!75, colbacktitle=liagreen, title=#1}

% Karte für die ganze Spaltenbreite (z. B. Mini-Kurs am Ende)
\newtcolorbox{liawide}[1]{
  enhanced,
  colback=white,
  colframe=liagreen!75,
  boxrule=0.3pt, arc=2pt,
  left=4pt,right=4pt,top=2pt,bottom=2pt,
  before skip=3pt, after skip=0pt,
  breakable=false,
  fonttitle=\sffamily\bfseries\footnotesize,
  coltitle=white, colbacktitle=liagreen, title=#1
}

% Info-Box ohne sidebyside (für Tools/Templates-Spalten unten)
\tcbset{
  infobase/.style={
    enhanced, colback=white,
    boxrule=0.3pt, arc=2pt,
    left=5pt,right=5pt,top=3pt,bottom=3pt,
    before skip=3pt, after skip=0pt,
    breakable=false,
    fonttitle=\sffamily\bfseries\footnotesize, coltitle=white
  }
}
\newtcolorbox{mdinfo}[1]{infobase, colframe=mdblue!75,    colbacktitle=mdblue,   title=#1}
\newtcolorbox{liainfo}[1]{infobase, colframe=liagreen!75, colbacktitle=liagreen, title=#1}

% Kleine Hilfsmacros für die "Wirkung"-Spalte
\newcommand{\prevtxt}[1]{{\small #1}}
\newcommand{\prevsm}[1]{{\footnotesize\color{muted} #1}}

\pagestyle{empty}
\begin{document}

% =====================================================================
% SEITE 1 — MARKDOWN
% =====================================================================
\noindent\begin{tcolorbox}[
  enhanced, frame hidden, sharp corners,
  colback=mdblue, coltext=white,
  left=12pt, right=8pt, top=4pt, bottom=4pt,
  before skip=0pt, after skip=4pt,
  width=\linewidth,
  borderline west={4pt}{-4pt}{mdblue!55}
]
\begin{minipage}[c]{0.66\linewidth}
  {\sffamily\fontsize{26}{28}\selectfont\bfseries Markdown}\hspace{0.4em}%
  {\sffamily\large\color{white!75} — die Basis}\\[2pt]
  {\sffamily\footnotesize\itshape\color{white!90} Was jeder Markdown-Editor versteht — Grundlage für alle LiaScript-Dokumente.}
\end{minipage}\hfill
\begin{minipage}[c]{0.20\linewidth}\raggedright
  {\sffamily\scriptsize\color{white!85}\faGraduationCap~~Hands-on Lab~$\cdot$~BiblioCon 2026}\\[2pt]
  {\sffamily\bfseries\normalsize Seite 1\,/\,2}\\[1pt]
  {\sffamily\scriptsize\color{white!75} Markdown-Grundlagen}
\end{minipage}\hfill
\begin{minipage}[c]{0.10\linewidth}\centering
  {\color{black}\colorbox{white}{\hspace{1pt}\qrcode[height=1.55cm,nolinks]{https://liascript.github.io/LiveEditor/}\hspace{1pt}}}\\[1pt]
  {\sffamily\scriptsize\color{white!90}\faEdit~Live-Editor}
\end{minipage}
\end{tcolorbox}

% --- Tag für Aufgabenbezug (klein, oben rechts in der Titelzeile) ---
\newcommand{\aufg}[1]{\hfill{\sffamily\scriptsize\color{white!85}\faTasks~A#1}}

\begin{multicols}{3}

% =====================================================================
% Block A — Textauszeichnung & Struktur (Aufgaben 2-4)
% =====================================================================

% --- A2a. Textauszeichnung ------------------------------------------
\begin{mdcard}{Textauszeichnung\aufg{2}}
\begin{lstlisting}[language=lia]
**fett**, *kursiv*,
~~durchgestrichen~~,
`Quellcode`
\end{lstlisting}
\tcblower
\prevtxt{\textbf{fett}, \textit{kursiv},
\sout{durchgestrichen}, \texttt{Quellcode}}
\end{mdcard}

% --- A2b. Absätze & Trennlinie --------------------------------------
\begin{mdcard}{Absätze \& Trennlinie\aufg{2}}
\begin{lstlisting}[language=lia]
Erster Absatz.

Zweiter Absatz nach
einer Leerzeile.

---
\end{lstlisting}
\tcblower
\prevtxt{Erster Absatz.\\[3pt]
Zweiter Absatz nach einer Leerzeile.\\[2pt]
\textcolor{rulegray}{\hrulefill}}
\end{mdcard}

% --- A3. Überschriften ----------------------------------------------
\begin{mdcard}{Überschriften\aufg{3}}
\begin{lstlisting}[language=lia]
# Kapitel
## Abschnitt
### Unterabschnitt
#### Detail
\end{lstlisting}
\tcblower
{\large\bfseries Kapitel}\par\vspace{1pt}
{\normalsize\bfseries Abschnitt}\par\vspace{1pt}
{\small\bfseries Unterabschnitt}\par\vspace{1pt}
{\footnotesize\bfseries Detail}
\end{mdcard}

% --- A4. Listen -----------------------------------------------------
\begin{mdcard}{Listen\aufg{4}}
\begin{lstlisting}[language=lia]
- erstens
- zweitens

  - eingerueckt

- drittens

1. eins
2. zwei
\end{lstlisting}
\tcblower
\prevtxt{%
$\bullet$ erstens\\
$\bullet$ zweitens\\[2pt]
\hspace*{1em}$\circ$ eingerückt\\[2pt]
$\bullet$ drittens\\[3pt]
1. eins\\
2. zwei\\[2pt]
{\scriptsize\color{muted}Untergeordnete Listen um \textbf{2 Zeichen} einrücken.}}
\end{mdcard}

% =====================================================================
% Block B — Tabellen, Code & Formeln (Aufgaben 5-6)
% =====================================================================

% --- A5. Tabellen ---------------------------------------------------
\begin{mdcard}{Tabellen\aufg{5}}
\begin{lstlisting}[language=lia]
| Operator | Wirkung      |
|:---------|:------------:|
| AND      | verkleinert  |
| OR       | vergroessert |
| NOT      | schliesst aus|
\end{lstlisting}
\tcblower
\prevtxt{%
\renewcommand{\arraystretch}{1.05}%
\begin{tabular}{|l|c|}\hline
\textbf{Operator} & \textbf{Wirkung}\\\hline
AND & verkleinert\\
OR  & vergrößert\\
NOT & schließt aus\\\hline
\end{tabular}\\[2pt]
{\scriptsize\color{muted}\texttt{:--} links, \texttt{:--:} zentriert, \texttt{--:} rechts}}
\end{mdcard}

% --- A6a. Inline-Code & Codeblock -----------------------------------
\begin{mdcard}{Inline-Code \& Codeblock\aufg{6}}
\begin{lstlisting}[language=lia]
Operatoren wie `AND`
und `OR` als Code.

```python
print("Hallo")
```
\end{lstlisting}
\tcblower
\prevtxt{Operatoren wie \texttt{\small AND}
und \texttt{\small OR} als Code.\\[3pt]
\colorbox{codebg}{\parbox{0.85\linewidth}{\ttfamily\scriptsize print("Hallo")}}}
\end{mdcard}

% --- A6b. Mathematik ------------------------------------------------
\begin{mdcard}{Mathematische Formeln\aufg{6}}
\begin{lstlisting}[language=lia]
Inline: $A \land B$

Block:
$$\neg A \lor (B \land C)$$
\end{lstlisting}
\tcblower
\prevtxt{Inline: $A \land B$\\[3pt]
Block:\\
$\displaystyle \neg A \lor (B \land C)$\\[3pt]
{\scriptsize\color{muted}\texttt{\textbackslash land}~$\land$,\,
\texttt{\textbackslash lor}~$\lor$,\,
\texttt{\textbackslash neg}~$\neg$}}
\end{mdcard}

% =====================================================================
% Block C — Links, Bilder, Hinweise (Aufgaben 7-8)
% =====================================================================

% --- A7a. Links ------------------------------------------------------
\begin{mdcard}{Links\aufg{7}}
\begin{lstlisting}[language=lia]
[Beschreibung](https://...)

<https://example.org>
\end{lstlisting}
\tcblower
\prevtxt{\textcolor{mdblue}{\underline{Beschreibung}}\\[2pt]
\textcolor{mdblue}{\underline{https://example.org}}\\[3pt]
{\scriptsize\color{muted}Sprechende Beschreibung erleichtert die Orientierung.}}
\end{mdcard}

% --- A7b. Bilder -----------------------------------------------------
\begin{mdcard}{Bilder\aufg{7}}
\begin{lstlisting}[language=lia]
![Alt-Text](pfad/bild.png)

![Venn-Diagramm AND/OR](
  https://.../venn.svg)
\end{lstlisting}
\tcblower
\prevtxt{\fbox{\rule{0pt}{1.2em}\hspace{1.2em}{\scriptsize\color{muted}Bild}\hspace{1.2em}}\\[2pt]
{\scriptsize\color{muted}\textit{Alt-Text in eckigen Klammern — wichtig für Barrierefreiheit.}}}
\end{mdcard}

% --- A8. Zitate & Hinweise -----------------------------------------
\begin{mdcard}{Zitate \& Hinweise\aufg{8}}
\begin{lstlisting}[language=lia]
> Ein einfaches Zitat.

> [!NOTE]    blau
> [!TIP]     gruen
> [!IMPORTANT] lila
> [!WARNING] orange
> [!CAUTION] rot
\end{lstlisting}
\tcblower
\prevtxt{%
\textcolor{rulegray}{\rule{1pt}{1.1em}}\hspace{4pt}\textit{Ein einfaches Zitat.}\\[2pt]
\textcolor{mdblue}{\rule{1pt}{1.0em}}\hspace{3pt}{\scriptsize\textbf{NOTE}\,/\,\textbf{TIP}\,/\,\textbf{IMPORTANT}}\\
\hspace*{8pt}{\scriptsize\textbf{WARNING}\,/\,\textbf{CAUTION}}\\
\hspace*{8pt}{\scriptsize Folgezeilen = Inhalt der Box.}}
\end{mdcard}

% =====================================================================
% Block D — Referenz-Karten (in den Aufgaben nicht zwingend gebraucht)
% =====================================================================

% --- R1. Aufgabenliste / Checkboxen ---------------------------------
\begin{mdcard}{Aufgabenliste}
\begin{lstlisting}[language=lia]
- [x] erledigt
- [ ] offen
- [ ] noch offen
\end{lstlisting}
\tcblower
\prevtxt{%
\faCheckSquare~erledigt\\
\faSquare~offen\\
\faSquare~noch offen}
\end{mdcard}

% --- R2. Fußnoten ---------------------------------------------------
\begin{mdcard}{Fußnoten}
\begin{lstlisting}[language=lia]
Wichtig[^1] zum Lernen.

[^1]: Quelle: ...
\end{lstlisting}
\tcblower
\prevtxt{Wichtig\textsuperscript{\textcolor{mdblue}{1}} zum Lernen.\\[3pt]
{\scriptsize \textsuperscript{1}~Quelle: \dots}}
\end{mdcard}

% --- R3. Definitionslisten ------------------------------------------
\begin{mdcard}{Definitionen}
\begin{lstlisting}[language=lia]
OER
: Open Educational
  Resources

AND
: Boole'scher Operator
  fuer Schnittmenge.
\end{lstlisting}
\tcblower
\prevtxt{\textbf{OER}\\
\hspace*{1em}Open Educational Resources\\[3pt]
\textbf{AND}\\
\hspace*{1em}Boole'scher Operator für Schnittmenge.}
\end{mdcard}

% --- R4. HTML & Sonderzeichen ---------------------------------------
\begin{mdcard}{HTML \& Sonderzeichen}
\begin{lstlisting}[language=lia]
<u>unterstrichen</u>
<br>  Zeilenumbruch
\* kein Listenpunkt
\end{lstlisting}
\tcblower
\prevtxt{\underline{unterstrichen}\\
{\scriptsize\color{muted}(Zeilenumbruch)}\\
*~kein Listenpunkt\\[2pt]
{\scriptsize\color{muted}Backslash escapt Markdown-Sonderzeichen.}}
\end{mdcard}

\end{multicols}

% --- Untere Zeile: Werkzeuge & Editoren (volle Breite) -------------
\noindent
\begin{mdinfo}{\faToolbox~~Werkzeuge \& Editoren für Markdown}
\footnotesize
\begin{minipage}[t]{0.24\linewidth}
\textbf{\color{mdblue}Lokale Editoren}\\[1pt]
$\bullet$~VS Code \textit{(+ Markdown All in One)}\\
$\bullet$~Obsidian, Typora, Marktext
\end{minipage}\hfill
\begin{minipage}[t]{0.22\linewidth}
\textbf{\color{mdblue}Online \& kollaborativ}\\[1pt]
$\bullet$~dillinger.io, stackedit.io\\
$\bullet$~hackmd.io \textit{(kollaborativ)}
\end{minipage}\hfill
\begin{minipage}[t]{0.24\linewidth}
\textbf{\color{mdblue}Konvertieren \& Qualität}\\[1pt]
$\bullet$~\texttt{pandoc} → HTML, PDF, DOCX, EPUB\\
$\bullet$~\texttt{markdownlint}, \texttt{prettier}
\end{minipage}\hfill
\begin{minipage}[t]{0.25\linewidth}
\textbf{\color{mdblue}Versionierung}\\[1pt]
$\bullet$~Git \& GitHub/GitLab rendern Markdown direkt\\
$\bullet$~Tabellen: \texttt{tablesgenerator.com/markdown}
\end{minipage}
\end{mdinfo}

\vfill
{\sffamily\scriptsize\color{muted}
\faInfoCircle~Live ausprobieren: \textbf{liascript.github.io/LiveEditor}\hfill
\faGithub~Quellcode: \textbf{github.com/LiaPlayground/Bibliocon2026}}

\newpage

% =====================================================================
% SEITE 2 — LIASCRIPT
% =====================================================================
\noindent\begin{tcolorbox}[
  enhanced, frame hidden, sharp corners,
  colback=liagreen, coltext=white,
  left=12pt, right=8pt, top=4pt, bottom=4pt,
  before skip=0pt, after skip=4pt,
  width=\linewidth,
  borderline west={4pt}{-4pt}{liagreen!55}
]
\begin{minipage}[c]{0.66\linewidth}
  {\sffamily\fontsize{26}{28}\selectfont\bfseries LiaScript}\hspace{0.4em}%
  {\sffamily\large\color{white!75} — die Erweiterungen}\\[2pt]
  {\sffamily\footnotesize\itshape\color{white!90} Alles, was Markdown allein nicht kann: Quiz, Animationen, ausführbarer Code, Sprache.}
\end{minipage}\hfill
\begin{minipage}[c]{0.20\linewidth}\raggedright
  {\sffamily\scriptsize\color{white!85}\faGraduationCap~~Hands-on Lab~$\cdot$~BiblioCon 2026}\\[2pt]
  {\sffamily\bfseries\normalsize Seite 2\,/\,2}\\[1pt]
  {\sffamily\scriptsize\color{white!75} Erweiterungen \& Interaktion}
\end{minipage}\hfill
\begin{minipage}[c]{0.10\linewidth}\centering
  {\color{black}\colorbox{white}{\hspace{1pt}\qrcode[height=1.55cm,nolinks]{https://github.com/LiaPlayground/Bibliocon2026}\hspace{1pt}}}\\[1pt]
  {\sffamily\scriptsize\color{white!90}\faGithub~Workshop-Repo}
\end{minipage}
\end{tcolorbox}

% --- Tag für Aufgabenbezug (grün, oben rechts im Karten-Titel) -----
\newcommand{\aufglia}[1]{\hfill{\sffamily\scriptsize\color{white!85}\faTasks~A#1}}

\begin{multicols}{3}

% =====================================================================
% Block A — Kursrahmen (Referenz, ohne direkten Aufgaben-Bezug)
% =====================================================================

% --- YAML-Header -----------------------------------------------------
\begin{liacard}{YAML-Header\aufglia{1}}
\begin{lstlisting}[language=lia]
<!--
author:   Maria Muster
language: de
narrator: Deutsch Female
import:   https://github.com/
          LiaTemplates/Pyodide
-->
\end{lstlisting}
\tcblower
\prevsm{Steht ganz oben im Dokument.\\
Definiert Autor:in, Sprache, Stimme
und externe Bausteine (\textit{Imports}).}
\end{liacard}

% =====================================================================
% Block B — LiaScript-Erweiterungen (Aufgaben 9-11)
% =====================================================================

% --- A9a. Medien abspielbar -----------------------------------------
\begin{liacard}{Medien abspielbar\aufglia{9}}
\begin{lstlisting}[language=lia]
!?[Titel](https://www.youtube
.com/watch?v=VIDEO_ID)

!?[Audio](datei.mp3)
\end{lstlisting}
\tcblower
\prevsm{\faYoutube~YouTube, Vimeo, MP4
und MP3 werden direkt eingebettet.\\[2pt]
\textbf{!?} statt \textbf{!} = abspielbar.}
\end{liacard}

% --- A9b. Eingebettete Ressource ------------------------------------
\begin{liacard}{Eingebettete Ressource\aufglia{9}}
\begin{lstlisting}[language=lia]
??[Katalog](https://kxp.k10plus.de)

??[Spezifikation](
  https://dublincore.org/...)
\end{lstlisting}
\tcblower
\prevsm{\faWindowMaximize~Erkennt automatisch
das Format (YouTube, PDF, GeoGebra,
oEmbed, \dots) -- iframe nur als
Fallback.\\[2pt]
\textbf{??} statt \textbf{?} = eingebettet
statt nur verlinkt.}
\end{liacard}

% --- A10. Animation / Click-Reveal ----------------------------------
\begin{liacard}{Animation \& Click-Reveal\aufglia{10}}
\begin{lstlisting}[language=lia]
Sofort sichtbar.

{{1}}
Erst nach einem Klick.

{{2}}
Nach dem zweiten Klick.
\end{lstlisting}
\tcblower
\prevsm{Sofort sichtbar.\\
{\color{rulegray}---}\\
\faMousePointer~nach 1. Klick: \textit{\dots}\\
\faMousePointer~nach 2. Klick: \textit{\dots}}
\end{liacard}

% --- A11. Sprecher-Text ---------------------------------------------
\begin{liacard}{Sprecher-Text (TTS)\aufglia{11}}
\begin{lstlisting}[language=lia]
## Folie
         --{{0}}--
Dieser Text wird vorgelesen,
sobald die Folie erscheint.
\end{lstlisting}
\tcblower
\prevsm{\faVolumeUp~Wird automatisch
gesprochen.\\
\texttt{\{\{0\}\}}~=~beim Öffnen,
\texttt{\{\{1\}\}}~=~nach 1.~Klick~\dots\\[2pt]
{\footnotesize Stimme im Header: \texttt{narrator:}}}
\end{liacard}

% =====================================================================
% Block C — Quizze (Aufgaben 12-13)
% =====================================================================

% --- A12. Freitext --------------------------------------------------
\begin{liacard}{Quiz: Freitext\aufglia{12}}
\begin{lstlisting}[language=lia]
Welcher Operator verkleinert
die Treffermenge?
[[AND]]
[[?]] Beide Begriffe muessen
      vorkommen.
\end{lstlisting}
\tcblower
\prevsm{Welcher Operator verkleinert die Treffermenge?\\[2pt]
\fbox{\strut\hspace{6em}}\\[3pt]
\faQuestionCircle~Hinweis (\texttt{[[?]]}, klappbar)\\[2pt]
{\footnotesize Mehrere gültige Antworten: weitere \texttt{[[\dots]]} anhängen.}}
\end{liacard}

% --- A13a. Single-Choice --------------------------------------------
\begin{liacard}{Quiz: Single-Choice\aufglia{13}}
\begin{lstlisting}[language=lia]
- [( )] falsch
- [(X)] richtig
- [[?]] Hinweis
*******************

Bravo, alles richtig!

*******************
\end{lstlisting}
\tcblower
\prevsm{$\circ$~falsch\\
$\bullet$~\textbf{richtig}\\
\faQuestionCircle~Hinweis (klappbar)\\[3pt]
{\footnotesize\color{liagreen}\faCheck~\textit{Bravo, alles richtig!}}\\[2pt]
{\footnotesize Lösungsblock nach \texttt{*****}\,---\,erscheint erst bei korrekter Antwort.}}
\end{liacard}

% --- A13b. Multiple-Choice ------------------------------------------
\begin{liacard}{Quiz: Multiple-Choice\aufglia{13}}
\begin{lstlisting}[language=lia]
- [[X]] vergroessert Treffer
- [[X]] verknuepft Synonyme
- [[ ]] schliesst Begriffe aus
- [[X]] ist Vereinigungsmenge
\end{lstlisting}
\tcblower
\prevsm{\faCheckSquare~vergrößert Treffer\\
\faCheckSquare~verknüpft Synonyme\\
\faSquare~schließt Begriffe aus\\
\faCheckSquare~ist Vereinigungsmenge\\[2pt]
{\footnotesize \texttt{[[X]]} = beliebig viele richtig.}}
\end{liacard}

% =====================================================================
% Block D — Diagramme, Zitation & Templates (Aufgaben 14-15)
% =====================================================================

% --- A14. Diagramme (Mermaid) ---------------------------------------
\begin{liacard}{Diagramme (Mermaid)\aufglia{14}}
\begin{lstlisting}[language=lia]
```mermaid @mermaid
graph TD
A --> B((OR))
B --> C[Bibliothek]
B --> D[Lesesaal]
```
\end{lstlisting}
\tcblower
\prevsm{\faProjectDiagram~Mermaid, ABC-Musik,
Plot, ASCII-Art --- alles inline.\\[2pt]
{\footnotesize \textbf{Header:} \texttt{import:}\\
\texttt{.../liaScript/mermaid\_template}}}
\end{liacard}

% --- A15. Zitation & BibTeX -----------------------------------------
\begin{liacard}{Zitation \& BibTeX\aufglia{15}}
\begin{lstlisting}[language=lia]
Siehe @cite{gantert2016}.

```bib
@book{gantert2016,
  author = {Klaus Gantert},
  title  = {Bibl. Grundwissen},
  year   = {2016}
}
```
\end{lstlisting}
\tcblower
\prevsm{Siehe \textcolor{liagreen}{[Gantert 2016]}.\\[3pt]
{\footnotesize BibTeX-Block am Ende des Dokuments —
LiaScript erzeugt automatisch das Literaturverzeichnis.}}
\end{liacard}

% --- Code ausführen (Pyodide) — Referenz aus Phase 1 ----------------
\begin{liacard}{Code ausführen}
\begin{lstlisting}[language=lia]
```python
zahlen = [1,2,3,4,5]
print(sum(zahlen))
```
@Pyodide.eval
\end{lstlisting}
\tcblower
\prevsm{\colorbox{codebg}{\parbox{0.85\linewidth}{\ttfamily\scriptsize zahlen = [1,2,3,4,5]\\print(sum(zahlen))}}\\[2pt]
\faPlay~\textbf{Run}~~\textcolor{liagreen}{\texttt{> 15}}\\[2pt]
{\footnotesize Setzt \texttt{Pyodide}-Import voraus.}}
\end{liacard}

\end{multicols}

% --- Untere Zeile: LiaScript-Templates (volle Breite) ---------------
\noindent
\begin{liainfo}{\faPuzzlePiece~~LiaScript-Templates — im Header per \texttt{import:} einbinden (Auswahl aus \texttt{github.com/LiaTemplates})}
\footnotesize
\begin{minipage}[t]{0.24\linewidth}
\textbf{\color{liagreen}\faTerminal~~Code ausführen}\\[1pt]
$\bullet$~\texttt{Pyodide}~–~Python\\
$\bullet$~\texttt{Skulpt}~–~Python (leichter)\\
$\bullet$~\texttt{Lua}~–~Lua\\
$\bullet$~\texttt{JSCPP}~–~C++\\
$\bullet$~\texttt{Tau-Prolog}~–~Prolog
\end{minipage}\hfill
\begin{minipage}[t]{0.24\linewidth}
\textbf{\color{liagreen}\faChartBar~~Visualisieren}\\[1pt]
$\bullet$~\texttt{mermaid\_template}~–~Diagramme\\
$\bullet$~\texttt{JSXGraph}~–~Mathematik\\
$\bullet$~\texttt{Vega}~–~Datenvisualisierung\\
$\bullet$~\texttt{Tikz-Jax}~–~Tikz-Grafiken\\
$\bullet$~\texttt{ABCjs}~–~Musiknotation
\end{minipage}\hfill
\begin{minipage}[t]{0.24\linewidth}
\textbf{\color{liagreen}\faDatabase~~Daten \& Bibliothek}\\[1pt]
$\bullet$~\texttt{citations}~–~BibTeX/CSL\\
$\bullet$~\texttt{Wikimedia}~–~Commons-Einbettung\\
$\bullet$~\texttt{Comunica}~–~SPARQL / Linked Data\\
$\bullet$~\texttt{AlaSQL}~/~\texttt{SQLite}~–~SQL\\
$\bullet$~\texttt{Random}~–~Quiz-Banken
\end{minipage}\hfill
\begin{minipage}[t]{0.24\linewidth}
\textbf{\color{liagreen}\faGithub~~Übersicht}\\[1pt]
$\bullet$~\texttt{github.com/LiaTemplates}\\
$\bullet$~\texttt{github.com/LiaTemplates/Index}\\
$\bullet$~\texttt{github.com/LiaPlayground}\\
$\bullet$~Live-Editor:\\
\hspace*{0.6em}\texttt{liascript.github.io/LiveEditor}
\end{minipage}
\end{liainfo}

\vspace{4pt}
{\sffamily\scriptsize\color{muted}
\faInfoCircle~Live-Editor: \textbf{liascript.github.io/LiveEditor}\hfill
\faBook~Doku: \textbf{liascript.github.io}\hfill
\faGithub~\textbf{github.com/LiaScript} \quad
CC-BY-SA 4.0 — Sebastian Zug}

\end{document}
